\section{The Physical Meaning of the Debye Length}

Plasma physics research has been strongly motivated by controlled fusion for more than 70 years, with contemporary developments focusing on magnetic confinement (tokamaks, stellarators) and inertial confinement (lasers, Z-pinches). NERS research emphasizes intense laser–plasma interactions (ZEUS), pulsed-power and high-power microwave sources, as well as low-temperature plasmas (GEC).


In a gaseous state, the average interparticle distance should be much greater than the Bohr radius ($n^{-\frac{1}{3}} \gg a_0$. This is because atoms and ions will behave mostly independently if further apart than the Bohr radius, which dictates the natural scale of bound electron orbitals. In an ionized state, the thermal energy must be greater than or equal to the ionization energy ($k_BT \geq E_i$).

\begin{prop}
  A plasma necessarily has a large number of electrons within a Debye sphere.
\end{prop}

\begin{eg}
  Tokamak Plasma
\end{eg}
\begin{explanation}
  In these plasmas, 
\end{explanation}


\section{Non-Dimensionalization of Equations of Scales}
For all the problems we've been considering, we start with Poisson's equation. So, let's take a moment to just look at Poisson's equation on its own.

\begin{eg}
  Ablation Plasma Ion Implantation (APII)
\end{eg}
\begin{explanation}
  For this example, we choose that
\end{explanation}

\begin{eg}
  Plasma (ref. Lecture 1)
\end{eg}
\begin{explanation}
  For this example, we choose that
\end{explanation}

\begin{eg}
  Diode
\end{eg}
\begin{explanation}
  For this example, we choose that
\end{explanation}

\begin{eg}
  Laser-Plasma Interaction
\end{eg}
\begin{explanation}
  For this example, we choose that
\end{explanation}
